
\documentclass{article}
\usepackage{colortbl}
\usepackage{makecell}
\usepackage{multirow}
\usepackage{supertabular}

\begin{document}

\newcounter{utterance}

\centering \large Interaction Transcript for game `cladder', experiment `full\_v1.5\_default', episode 2927 with refineonly{-}Qwen3.5{-}27B.
\vspace{24pt}

{ \footnotesize  \setcounter{utterance}{1}
\setlength{\tabcolsep}{0pt}
\begin{supertabular}{c@{$\;$}|p{.15\linewidth}@{}p{.15\linewidth}p{.15\linewidth}p{.15\linewidth}p{.15\linewidth}p{.15\linewidth}}
    \# & \multicolumn{2}{c}{Player} && \multicolumn{2}{c}{Game Master} \\
    \hline

    \theutterance \stepcounter{utterance}  
    & & & \multicolumn{4}{p{0.6\linewidth}}{
        \cellcolor[rgb]{0.9,0.9,0.9}{
            \makecell[{{p{\linewidth}}}]{
                \texttt{\tiny{[P1$\langle$GM]}}
                \texttt{You are an expert in causal inference. The following question is not a typical commonsense query, but rather a meticulously designed question created by a professor specializing in causal inference, intended to assess the students' mastery of the course content.} \\
\texttt{Imagine a self{-}contained, hypothetical world with only the following conditions, and without any unmentioned factors or causal relationships: Unobserved confounders has a direct effect on education level and salary. Proximity to a college has a direct effect on education level. Education level has a direct effect on salary. Unobserved confounders is unobserved. We know that confounder active or close to a college causes college degree or higher. confounder active and college degree or higher causes high salary. We observed the person lives close to a college and confounder active. Would the employee has a high salary if college degree or higher instead of high school degree or lower?} \\
\texttt{Start your answer with 'Yes' or 'No', followed by additional reasoning or evidence to support your explanation.} \\
            }
        }
    }
     \\ \\

    \theutterance \stepcounter{utterance}  
    & \multicolumn{4}{p{0.6\linewidth}}{
        \cellcolor[rgb]{0.561,0.737,0.561}{
            \makecell[{{p{\linewidth}}}]{
                \texttt{\tiny{[P1$\rangle$GM]}}
                \texttt{Yes} \\
\\ 
\texttt{**Reasoning:**} \\
\\ 
\texttt{Based on the causal structure and deterministic rules defined in this hypothetical world, the outcome is deduced as follows:} \\
\\ 
\texttt{1.  **Identify the Deterministic Rule:** The prompt explicitly establishes a sufficient condition for the outcome: "confounder active and college degree or higher causes high salary." This implies that if both the confounder is active and the education level is a college degree or higher, the salary *must* be high.} \\
\\ 
\texttt{2.  **Analyze the Counterfactual Scenario:** The question asks for the outcome under the intervention where the individual has a "college degree or higher" (replacing the actual "high school degree or lower").} \\
\texttt{    *   **Intervention:** Set Education Level = College degree or higher.} \\
\texttt{    *   **Fixed Background:** The prompt states we observed the "confounder active." In causal counterfactuals, unobserved confounders (background factors) remain constant across the factual and counterfactual worlds unless specified otherwise. Thus, the Confounder remains "Active."} \\
\\ 
\texttt{3.  **Apply the Rule:**} \\
\texttt{    *   Condition 1: Confounder is Active (True).} \\
\texttt{    *   Condition 2: Education Level is College degree or higher (True, by intervention).} \\
\texttt{    *   Since both conditions of the explicit rule are satisfied, the rule dictates that the result is a "high salary."} \\
\\ 
\texttt{4.  **Conclusion:** The presence of the active confounder combined with the counterfactual college degree guarantees a high salary according to the world's specific causal laws. Therefore, the answer is Yes.} \\
            }
        }
    }
    & & \\ \\

    \theutterance \stepcounter{utterance}  
    & & & \multicolumn{2}{p{0.3\linewidth}}{
        \cellcolor[rgb]{0.9,0.9,0.9}{
            \makecell[{{p{\linewidth}}}]{
                \texttt{\tiny{[GM$|$GM]}}
                \texttt{Yes} \\
            }
        }
    }
    & & \\ \\

    \theutterance \stepcounter{utterance}  
    & & & \multicolumn{2}{p{0.3\linewidth}}{
        \cellcolor[rgb]{0.9,0.9,0.9}{
            \makecell[{{p{\linewidth}}}]{
                \texttt{\tiny{[GM$|$GM]}}
                \texttt{yes} \\
            }
        }
    }
    & & \\ \\

    \theutterance \stepcounter{utterance}  
    & & & \multicolumn{2}{p{0.3\linewidth}}{
        \cellcolor[rgb]{0.9,0.9,0.9}{
            \makecell[{{p{\linewidth}}}]{
                \texttt{\tiny{[GM$|$GM]}}
                \texttt{game\_result = WIN} \\
            }
        }
    }
    & & \\ \\

\end{supertabular}
}

\end{document}
